\fichatitulo{36 · REPRESENTAÇÃO}{NAACL · 2019}{BERT}
{\small Devlin et al. \textit{{BERT}: Pre-training of Deep Bidirectional Transformers for Language Understanding}. NAACL · 2019. \href{https://aclanthology.org/N19-1423.pdf}{Fonte consultada}.\par}

\campo{Problema e objetivo}
Como aprender representações contextualizadas bidirecionais transferíveis?

\campo{Principal contribuição · síntese interpretativa}
Propõe pré-treinamento bidirecional profundo, seguido de adaptação para tarefas de linguagem.

\campo{Método / natureza da evidência}
Combina modelagem de linguagem mascarada e previsão da próxima sentença; ajusta o modelo para tarefas supervisionadas.

\campo{Resultado ou argumento principal}
Relata melhorias em onze tarefas de processamento de linguagem natural.

\campo{Excertos-chave · seleção literal}
\begin{itemize}
\excerto{1}{left and right context in all layers}{Resumo.}
\excerto{2}{next sentence prediction}{Introdução.}
\end{itemize}

\campo{Limitações e análise crítica}
Análise da ficha: BERT não é um gerador autorregressivo completo nem um recuperador de passagens pronto. Usá-lo em busca requer definir representação, treinamento e função de similaridade.

\campo{Aproveitamento na dissertação · proposta}
Fundamentação: distinguir codificadores, geradores e adaptação para recuperação.

\registro{Fonte consultada nesta preparação: Resumo e introdução. Chave: \texttt{devlinEtAl2019bert}.}
